\documentclass[book.tex]{subfiles}

\begin{document}
\chapter{Creation and Annihilation Operators}

\label{chap:creation_annihilation}
\noindent\rule{11cm}{0.4pt} \\
\textbf{Prerequisites:} \autoref{chap:quantum_mechanics}, \autoref{chap:quantum_harmonic} \\
\textbf{Difficulty Level:} ** \\
\noindent\rule{11cm}{0.4pt}
\vspace{5mm}

In this chapter, we introduce a powerful mathematical technique that arises repeatedly in quantum systems: creation and annihilation operators. An annihilation operator is denoted $\hat{a}$ lowers the number of particles in a state by 1. A creation operator $\hat{a}^{\dagger}$ raises the number of particles by 1. Let's derive these operators in the context of the quantum harmonic oscillator we studied in \autoref{chap:quantum_harmonic}.

\begin{equation}
    \left (-\frac{\hbar^2}{2m}\frac{d^2}{dx^2} + \frac{1}{2}m\omega^2 x^2 \right ) \psi = E \psi
\end{equation}
Making a change of coordinates
\begin{equation}
    x = \sqrt{\frac{\hbar}{m\omega}}q
\end{equation}
Note that this means
\begin{align}
    dx &= \sqrt{\frac{\hbar}{m\omega}}dq \\
    \frac{dq}{dx} &= \sqrt{\frac{m\omega}{\hbar}} 
\end{align}
Plugging this into the change of variables, we have
\begin{align}
    \frac{d^2}{dx^2} &= \left (\frac{dq}{dx} \right)^2 \frac{d}{dq^2} \\ 
    \frac{d^2}{dx^2} &= \frac{m\omega}{\hbar} \frac{d}{dq^2}    
\end{align}
Substituting into the equation above, we get 
\begin{equation}
    \frac{\hbar \omega}{2} \left ( - \frac{d^2}{dq^2} + q^2 \right ) \psi(q) = E \psi(q) 
\end{equation}
The following factorization trick proves useful:
\begin{equation}
    \left (-\frac{d}{dq^2} + q^2 \right) = \left (-\frac{d}{dq} + q\right )\left (\frac{d}{dq} + q \right) + \frac{d}{dq} q - q \frac{d}{dq}
\end{equation}
Let's apply the second part of this to an arbitrary function $f$
\begin{align}
    \left (\frac{d}{dq} q - q \frac{d}{dq} \right ) f &= \frac{d}{dq}(qf(q)) - q\left(\frac{df(q)}{dq} \right )\\
    &= f(q) + q\frac{df(q)}{dq} - q\frac{df(q)}{dq} \\
    &= f(q)
\end{align}
Which implies that 
\begin{equation}
    \frac{d}{dq} q - q\frac{d}{dq} = 1
\end{equation}
It follows that
\begin{equation}
    -\frac{d^2}{dq^2} + q^2 = \left (-\frac{d}{dq} + q \right ) \left (\frac{d}{dq} + q \right )  + 1
\end{equation}
Now define
\begin{equation}
    \hat{a}^{\dagger} = \frac{1}{\sqrt{2}}\left ( -\frac{d}{dq} + q \right )
\end{equation}
and
\begin{equation}
    \hat{a}^ = \frac{1}{\sqrt{2}} \left ( \frac{d}{dq} + q \right )
\end{equation}
Then we can rewrite Schrodinger's equation as

\begin{equation}
    \hbar \omega (\hat{a}^{\dagger} \hat{a} + \frac{1}{2}) \psi = E \psi
\end{equation}


%\textcolor{red}{TODO: Add a discussion of the commutators of the creation and annihilation operators}

%\textcolor{red}{TODO: Complete discussion of how to use creation/annihilation operators to solve Schrodinger's equation}

\end{document}
